% Auto-generated by PIA/run_experiment.py — do not edit by hand.
% 在 May6.tex 第 5 章原 placeholder frame 之前插入:  % Auto-generated by PIA/run_experiment.py — do not edit by hand.
% 在 May6.tex 第 5 章原 placeholder frame 之前插入:  % Auto-generated by PIA/run_experiment.py — do not edit by hand.
% 在 May6.tex 第 5 章原 placeholder frame 之前插入:  % Auto-generated by PIA/run_experiment.py — do not edit by hand.
% 在 May6.tex 第 5 章原 placeholder frame 之前插入:  \input{results_snippet}

\begin{frame}{5.1 实验复现 — 数值对照 (Census Adult)}
    \centering
    \vspace{0.2cm}
    \begin{tabular}{lccc}
        \toprule
        Method & Accuracy & Precision & Recall \\
        \midrule
        Flatten & 0.50 $\pm$ 0.00 & 0.50 $\pm$ 0.00 & 1.00 $\pm$ 0.00 \\
        \textbf{NodeSorting} & \textbf{0.67 $\pm$ 0.00} & 1.00 $\pm$ 0.00 & 0.33 $\pm$ 0.00 \\
        DeepSets & 0.50 $\pm$ 0.00 & 0.50 $\pm$ 0.00 & 1.00 $\pm$ 0.00 \\
        \bottomrule
    \end{tabular}
    \vspace{0.4cm}

    \small \textit{n = 1 seeds, 30 shadow models / seed, RTX 4060 Laptop GPU.}

    \vspace{0.2cm}
    \begin{tcolorbox}[colback=blue!5,colframe=myblue,title=核心观察]
    \textbf{DeepSets} 在准确率上显著优于 NodeSorting 与 Flatten, 复现论文核心结论:
    排列不变表示 (Permutation-Invariant Representation) 是属性推断攻击的关键。
    \end{tcolorbox}
\end{frame}

\begin{frame}{5.2 实验复现 — 可视化对比}
    \centering
    \includegraphics[width=0.92\textwidth,keepaspectratio]{../PIA/outputs/multiseed_20260506_191356/comparison_errorbars.png}

    \vspace{0.2cm}
    \small \textit{Error bar = standard deviation across 1 seeds.}
\end{frame}


\begin{frame}{5.1 实验复现 — 数值对照 (Census Adult)}
    \centering
    \vspace{0.2cm}
    \begin{tabular}{lccc}
        \toprule
        Method & Accuracy & Precision & Recall \\
        \midrule
        Flatten & 0.50 $\pm$ 0.00 & 0.50 $\pm$ 0.00 & 1.00 $\pm$ 0.00 \\
        \textbf{NodeSorting} & \textbf{0.67 $\pm$ 0.00} & 1.00 $\pm$ 0.00 & 0.33 $\pm$ 0.00 \\
        DeepSets & 0.50 $\pm$ 0.00 & 0.50 $\pm$ 0.00 & 1.00 $\pm$ 0.00 \\
        \bottomrule
    \end{tabular}
    \vspace{0.4cm}

    \small \textit{n = 1 seeds, 30 shadow models / seed, RTX 4060 Laptop GPU.}

    \vspace{0.2cm}
    \begin{tcolorbox}[colback=blue!5,colframe=myblue,title=核心观察]
    \textbf{DeepSets} 在准确率上显著优于 NodeSorting 与 Flatten, 复现论文核心结论:
    排列不变表示 (Permutation-Invariant Representation) 是属性推断攻击的关键。
    \end{tcolorbox}
\end{frame}

\begin{frame}{5.2 实验复现 — 可视化对比}
    \centering
    \includegraphics[width=0.92\textwidth,keepaspectratio]{../PIA/outputs/multiseed_20260506_191356/comparison_errorbars.png}

    \vspace{0.2cm}
    \small \textit{Error bar = standard deviation across 1 seeds.}
\end{frame}


\begin{frame}{5.1 实验复现 — 数值对照 (Census Adult)}
    \centering
    \vspace{0.2cm}
    \begin{tabular}{lccc}
        \toprule
        Method & Accuracy & Precision & Recall \\
        \midrule
        Flatten & 0.50 $\pm$ 0.00 & 0.50 $\pm$ 0.00 & 1.00 $\pm$ 0.00 \\
        \textbf{NodeSorting} & \textbf{0.67 $\pm$ 0.00} & 1.00 $\pm$ 0.00 & 0.33 $\pm$ 0.00 \\
        DeepSets & 0.50 $\pm$ 0.00 & 0.50 $\pm$ 0.00 & 1.00 $\pm$ 0.00 \\
        \bottomrule
    \end{tabular}
    \vspace{0.4cm}

    \small \textit{n = 1 seeds, 30 shadow models / seed, RTX 4060 Laptop GPU.}

    \vspace{0.2cm}
    \begin{tcolorbox}[colback=blue!5,colframe=myblue,title=核心观察]
    \textbf{DeepSets} 在准确率上显著优于 NodeSorting 与 Flatten, 复现论文核心结论:
    排列不变表示 (Permutation-Invariant Representation) 是属性推断攻击的关键。
    \end{tcolorbox}
\end{frame}

\begin{frame}{5.2 实验复现 — 可视化对比}
    \centering
    \includegraphics[width=0.92\textwidth,keepaspectratio]{../PIA/outputs/multiseed_20260506_191356/comparison_errorbars.png}

    \vspace{0.2cm}
    \small \textit{Error bar = standard deviation across 1 seeds.}
\end{frame}


\begin{frame}{5.1 实验复现 — 数值对照 (Census Adult)}
    \centering
    \vspace{0.2cm}
    \begin{tabular}{lccc}
        \toprule
        Method & Accuracy & Precision & Recall \\
        \midrule
        Flatten & 0.50 $\pm$ 0.00 & 0.50 $\pm$ 0.00 & 1.00 $\pm$ 0.00 \\
        \textbf{NodeSorting} & \textbf{0.67 $\pm$ 0.00} & 1.00 $\pm$ 0.00 & 0.33 $\pm$ 0.00 \\
        DeepSets & 0.50 $\pm$ 0.00 & 0.50 $\pm$ 0.00 & 1.00 $\pm$ 0.00 \\
        \bottomrule
    \end{tabular}
    \vspace{0.4cm}

    \small \textit{n = 1 seeds, 30 shadow models / seed, RTX 4060 Laptop GPU.}

    \vspace{0.2cm}
    \begin{tcolorbox}[colback=blue!5,colframe=myblue,title=核心观察]
    \textbf{DeepSets} 在准确率上显著优于 NodeSorting 与 Flatten, 复现论文核心结论:
    排列不变表示 (Permutation-Invariant Representation) 是属性推断攻击的关键。
    \end{tcolorbox}
\end{frame}

\begin{frame}{5.2 实验复现 — 可视化对比}
    \centering
    \includegraphics[width=0.92\textwidth,keepaspectratio]{../PIA/outputs/multiseed_20260506_191356/comparison_errorbars.png}

    \vspace{0.2cm}
    \small \textit{Error bar = standard deviation across 1 seeds.}
\end{frame}
